% chapters/07_fazit.tex

\chapter{Fazit und Ausblick}
\label{chap:fazit}

\section{Fazit}
\label{sec:fazit}

Diese Projektarbeit ist der Frage nachgegangen, ob und wie sich das
Vision-Language-Action-Modell NVIDIA GR00T~N1.6 auf dem Unitree~G1 mit DEX3-Hand
für eine Block-Stacking-Aufgabe feinabstimmen lässt -- und welche Aussagekraft die
anschließende Simulationsevaluation besitzt. Die zentralen Erkenntnisse lassen
sich wie folgt zusammenfassen:

\begin{itemize}
  \item \textbf{Reproduzierbare Infrastruktur.} Es wurde eine containerisierte,
    rein über Umgebungsvariablen konfigurierte Pipeline aufgebaut, die identisch
    lokal (Docker) und auf dem KISSKI-Cluster (Apptainer/\ac{slurm}) läuft und
    auch auf vier GPUs skaliert. Sie macht das Training reproduzierbar und
    plattformunabhängig.
  \item \textbf{Erfolgreiches Training.} Das selektive Fine-Tuning (eingefrorenes
    \ac{vlm}, trainierter Flow-Matching-Aktionskopf) konvergierte über
    175\,000~Schritte sauber und ohne Instabilitäten; der Loss plateaut
    ab~$\approx$~100\,000~Schritten.
  \item \textbf{Aufgedeckter Domain-Gap als Kernbefund.} Trotz guter Open-Loop-
    Armvorhersagen erreichte die Politik im geschlossenen Regelkreis
    $0/20$~Erfolge und \enquote{fror ein}. Eine Dreifach-Diagnose
    (Open-Loop, Replay, Closed-Loop) isolierte die \emph{Beobachtung} als
    einzige Differenz, und eine SigLIP-basierte Messung quantifizierte den
    visuellen Domain-Gap -- mit der handnahen Kamera
    (\texttt{cam\_left\_wrist}, Distanz~0,427) als kritischem Blocker.
  \item \textbf{Einordnung statt Fehlinterpretation.} Literaturgestützt
    (SIMPLER~\cite{simpler2024}) wurde gezeigt, dass $0\,\%$ Closed-Loop-Erfolg
    für ein auf Realbildern trainiertes Modell mit eingefrorenem Encoder in einer
    nicht-photorealistischen Sim das \emph{erwartete} Ergebnis ist -- nicht ein
    Urteil über die Modellqualität. Die offizielle, domain-gap-freie Metrik
    (Open-Loop-\ac{mse}) deutet darauf hin, dass das Modell die Aufgabe im Prinzip
    gelernt hat.
  \item \textbf{Extern validierte Pipeline, eingegrenzter Domain-Gap.} Die
    Inferenz-Kette reproduzierte einen unabhängigen, publizierten Benchmark
    (RoboCasa GR-1: gemessen 47,7\,\% vs.\ erwartet 47,8\,\%,
    \cref{sec:robocasa}) und ist damit als Fehlerquelle ausgeschlossen.
    Kamera-Neukalibrierung und Albedo-Korrekturen senkten den mittleren
    Real$\to$Sim-Abstand anschließend unter die Real$\to$Real-Referenz
    (\cref{sec:neukalibrierung}); als dominante Gap-Ursache erwies sich die
    Materialfarbe (Albedo), nicht die Beleuchtung.
  \item \textbf{Verhaltensbeleg statt bloßer Bildähnlichkeit.} Das
    \texttt{span}-Gate (\cref{sec:span-gate}) trennte die beiden verbliebenen
    Erklärungen: Dieselbe Politik kommandiert auf \emph{echten} Bildern die
    volle Greifbewegung (Verhältnis~1,00), in der Simulation nur ein Fünftel.
    Die Politik hat das Greifen also gelernt und ruft es nur bei synthetischer
    Beobachtung nicht ab.
  \item \textbf{Checkpoint-Auswahl war messbar falsch.} Der dritte Lauf, der
    erste mit echter Testmenge, zeigte eine U-Kurve des Validierungsfehlers:
    Der letzte Checkpoint ist um 25\,\% schlechter als der beste
    (\cref{sec:lauf3}). Overfitting ist damit im Projekt erstmals gemessen
    statt vermutet.
  \item \textbf{RL-Pipeline implementiert, Lernwirkung offen.} Das als
    Konzept gestartete RL-Fine-Tuning wurde als FPO-Trainer umgesetzt und
    absolvierte eine vollständige Trainingsiteration auf institutseigener
    Blackwell-Hardware (\cref{sec:rl-stand}). Der zwischenzeitlich vermutete
    Greif-Physik-Blocker ist ausgeräumt (\cref{sec:grasp-befund}); vor einem
    produktiven Lauf steht nun das Co-Training (\cref{sec:cotrain}).
  \item \textbf{Ein widerrufener Fix als methodische Warnung.} Eine früher als
    Korrektur verbuchte Vorzeichen-Spiegelung der Fingergelenke erwies sich
    zuletzt als selbst fehlerhaft; sie ließ die Hand in allen bisherigen
    Sim-Läufen unvollständig schließen (\cref{sec:vorzeichen-widerruf}). Die
    simulationsseitigen Greifzahlen dieser Arbeit stehen dadurch unter einem
    ausdrücklichen Vorbehalt (\cref{sec:vorbehalt-greifzahlen}).
\end{itemize}

Der wichtigste Beitrag dieser Arbeit ist somit weniger eine hohe Erfolgsrate als
die \emph{methodisch saubere Diagnose}, warum die naheliegende Sim-Evaluation für
dieses Modell die falsche Messgröße ist -- ein Ergebnis, das in der konkreten
Empfehlung mündet, Modellqualität primär über Open-Loop-Eval und reale
Roboter-Rollouts zu belegen.

\section{Ausblick}
\label{sec:ausblick}

\paragraph{Messungen mit korrigierter Geometrie wiederholen.}
Der unmittelbarste nächste Schritt folgt aus dem Widerruf des Vorzeichen-Fixes
(\cref{sec:vorzeichen-widerruf}): Die simulationsseitigen Greifgrößen --
Erfolgsrate, Fingerspanne, Meilenstein-Leiter -- sind mit korrigierter
Hand- und Basisgeometrie neu zu erheben. Erst danach existiert ein belastbarer
Nullpunkt, gegen den sich Co-Training und RL überhaupt messen lassen, und erst
danach lässt sich beziffern, welcher Anteil des Sim-Versagens tatsächlich auf
den visuellen Domain-Gap entfällt (\cref{sec:vorbehalt-greifzahlen}).

\paragraph{Co-Training durchführen.}
Die Werkzeugkette für das Training auf echten \emph{und} gerenderten Bildern
steht vollständig (\cref{sec:cotrain}); der Lauf selbst steht aus. Er ist der
inhaltlich nächste Schritt, weil er als einziger die in
\cref{sec:frozen-vs-tuned} formulierte Bedingung erfüllt -- der Vision-Encoder
sieht während des Trainings tatsächlich Simulationsbilder. Der gerenderte
Datensatz ist dafür mit der korrigierten Geometrie neu zu erzeugen.

\paragraph{Domain-Gap weiter schließen.}
Visual Matching hat sich als wirksam erwiesen: Kamera-Neukalibrierung und
Albedo-Korrekturen senkten den mittleren Gap unter die Real$\to$Real-Referenz
(\cref{sec:neukalibrierung}). Für die verbleibende kritische Handgelenk-Kamera
ist der nächste Hebel ein Mittrainieren des Vision-Encoders
(\texttt{TUNE\_VISUAL}) -- allerdings nur \emph{mit} Simulations- oder
\acs{dr}-variierten Bildern: Auf reinen Realdaten schließt es den Gap nicht,
sondern verschlechterte das Verhalten bis zum Politik-Kollaps, wie der zweite
Lauf empirisch zeigte (\cref{sec:vision-lauf,sec:frozen-vs-tuned}). Genau diese
Bedingung erfüllt das Co-Training aus \cref{sec:cotrain}, dessen Werkzeugkette
bereitsteht.

\paragraph{Reinforcement-Learning-Fine-Tuning produktiv machen.}
Wie in \cref{sec:rl-weiterfuehrend} ausgeführt, adressiert RL die \ac{bc}-Schwäche
an der Wurzel und löst den Domain-Gap konzeptionell auf, da in derselben Sim
trainiert und evaluiert wird. Die ursprünglich kritische GPU-Frage ist durch den
institutseigenen Blackwell-Server beantwortet, die FPO-Pipeline implementiert
und end-to-end validiert (\cref{sec:rl-stand}). Nach dem Co-Training und den
wiederholten Basismessungen stehen der erste produktive Lauf, der Nachweis der
Lernwirkung und das Hyperparameter-Tuning an.

\paragraph{Lokomotion.}
Mit der entkoppelten Architektur (\cref{sec:loco-pfade}) lässt sich der bislang
fixierte Roboter zum Laufen bringen, ohne neue Manipulationsdaten zu erheben --
empfohlen als gestufter Weg von einer reinen Sim-Lauf-Demo über die entkoppelte
Integration bis hin, langfristig, zur unified Whole-Body-\ac{vla}.

\paragraph{Sim-to-Real.}
Der entscheidende, noch ausstehende Test ist der Einsatz auf der physischen
G1-Hardware. Da das Closed-Loop-Versagen ein sim-spezifisches Wahrnehmungsproblem
ist, ist plausibel, dass das Modell auf realen Kamerabildern -- die der
Trainingsverteilung entsprechen -- deutlich besser abschneidet. Der reale Roboter
bleibt der Goldstandard und der nächste große Meilenstein des Projekts.
